% ---------------------------------------------------------------------------
% The reverse direction in one page: this version's division of labour is the opposite
% of main_short.tex's, which carries these results with their proofs in full. Stated
% here because the asymmetry is what makes the forward direction's cheapness a finding
% rather than a reassurance, and because the decision procedure is the one place in
% this version where the course's machinery does the work.
% ---------------------------------------------------------------------------

\section{The Other Direction, Stated}
\label{s2:sec:reverse}

Everything above is one direction. The reverse direction is where language design does
bite, and it bites in a form that admits proof rather than measurement, so this section
states the results and leaves the proofs to the complete paper.

Both languages filter on tag values and both do it with a pattern language, but not the
same one. Datadog's is a \emph{glob}: a string $g=s_0*s_1*\cdots*s_n$ over
$\Sigma\cup\{*\}$ with literal segments $s_i$, denoting
$L(g)=\{\,s_0w_1s_1\cdots w_ns_n : w_i\in\Sigma^{*}\,\}$. Datadog documents prefix, suffix
and substring wildcards and no other construct~\citep{ddfilter} --- no single-character
wildcard, no alternation inside a value, no character class, and no escape making a literal
\texttt{*} --- which is the premise of every negative result here, and which was read from
the documentation until we put all three constructs to a real account, whose query parser
rejected all three. PromQL's is a regular expression, and replacing every \texttt{*} with
\texttt{.*} and escaping everything else gives one denoting the same language under
PromQL's whole-value anchoring, which is why no gap code in Table~\ref{s2:tab:gaps}
mentions globs on tag values. The reverse direction is the interesting one, and one
structural fact decides most of it: for a glob with no star $L(g)$ is a single word, and
for a glob with at least one star $L(g)$ contains a word of every length
$\ell\ge\sum_i|s_i|$ --- so a glob language is a singleton or it is unbounded in length.

\begin{corollary}[What globs cannot say]\label{s2:cor:no-glob}
Assume $|\Sigma|\ge 2$. No glob denotes a finite language with two or more words --- in
particular none denotes $\Sigma$, so there is no single-character wildcard, and none
denotes an alternation of distinct literals. No glob denotes $\Sigma^{+}$, the non-empty
values. For $\emptyset\neq A\subsetneq\Sigma$, no glob denotes $A^{+}$, so a star cannot be
restricted to part of the alphabet. And if $L$ is \emph{infinite} and some $\sigma\in\Sigma$
occurs in no word of $L$, then $L$ is not a union of finitely many glob languages.
\end{corollary}

\noindent
The first two parts follow from the dichotomy alone; the third needs what a star admits,
which is every character of $\Sigma$. The last is the only result here that holds against
an \emph{arbitrary} union, with no bound on how many branches a translator may spend, and
it is what separates ten of the refusals below from the rest. ``Infinite'' there is a
hypothesis and not a flourish: a finite language omits characters and is a union of globs,
one per word.

\begin{proposition}[Glob-expressibility is decidable]\label{s2:prop:decidable}
Let $r$ be a regular expression with $L(r)\neq\emptyset$, read over one fixed non-empty
$\Sigma$. It is decidable whether some glob $g$ has $L(g)=L(r)$, and a witness is produced
when one exists.
\end{proposition}

\noindent
One observation about shortest words settles it. If $L(r)=L(g)$ for $g=s_0*\cdots*s_n$,
taking every $w_i=\varepsilon$ gives $s_0s_1\cdots s_n\in L(g)$, and every other word of
$L(g)$ is strictly longer --- so that word is the shortest of $L(g)$ \emph{and the only
word of that length}. Build a DFA for $r$ by Rich's regular-expression-to-FSM construction
(Thm 6.1) and the subset construction (Thm 5.3), and find its shortest accepted word $w$,
with the number of accepting paths of that length, by breadth-first search. If there is no
unique shortest word, no glob exists; otherwise $g$ is $w$ with stars inserted at some
subset of its $|w|+1$ gaps --- at most $2^{|w|+1}$ candidates, each settled by deciding
equivalence of two regular languages (Rich Thm 9.5). The search is finite and exhaustive,
so it answers either way.

\paragraph{What the corpus contains, and the asymmetry.}
Of the $1{,}041$ concrete matcher occurrences in the PromQL half of the corpus --- the
eighth of the $8{,}587$ whose pattern is not a template variable --- a Datadog scope
reaches $91.9\%$, but two thirds of that only as a union of globs, running to $41$ branches
for the commonest case and $206$ at the widest. Of the $84$ refusals, $74$ rest on the
first three parts of Corollary~\ref{s2:cor:no-glob} together with a branch bound we choose
and state, and $10$ on its last part, which needs no bound at all. The procedure of
Proposition~\ref{s2:prop:decidable}, run over all $184$ distinct concrete patterns as a
cross-check, agrees with the syntax-directed proposer on every one.

Put the two directions beside each other. Forward, every candidate for inexpressibility was
removed by writing the code or by exhibiting the expression, bar $14$ occurrences left
undecided; backward, the residue divides into refusals resting on a bound we chose and
refusals resting on a theorem. The two halves are not symmetric in kind, and that is the
point: ``work remains'' is a statement about this translator, and
Remark~\ref{s2:rem:refusal} says such a statement is weak evidence about a language, while
``no glob denotes this language'' is a statement about globs that no amount of engineering
changes. So the lock-in is real and not where the question assumed it would be --- not that
a proprietary language says things an open one cannot, but that the proprietary filter
language is weaker, so the migration is one-way. Leaving costs work; returning costs
expressiveness, which is the testable form of the structural point
\S\ref{s2:sec:background} takes from the switching-cost literature: the terms of return
differ from the terms of entry. Two things limit how far that carries. The reverse direction
has no translator behind it, so its numbers describe what a converter could achieve rather
than what one does; and the ten unconditional refusals are ten, so a different corpus could
contain none while the theorems stood.
